\partsection{F}{Finite tests and the information retained by higher morphisms}\label{sec:F}

Let $\mathcal A_{\mathrm{cat}}\subset\Spc$ be a full, essentially small
$\infty$-subcategory containing one representative of every weak homotopy
type $BA$ with $A$ a category having finitely many objects and arrows.
Let $\mathcal A_{\mathrm{fin}}\subset\Spc$ be the full subcategory of finite
CW homotopy types. We choose compatible small skeleta when forming
presheaf categories. The ordinary category called $\Hot_0$ in Chapter~V
is $h\mathcal A_{\mathrm{cat}}$. Neither the word ``finite'' nor a
statement about $\Ind$ will be allowed to identify these two subcategories.
All the mapping spaces in this section are derived mapping spaces of
unpointed spaces; brackets denote their components.

\begin{claim}{F.1}{Proved}{The two relevant finiteness conditions}
Every finite CW homotopy type is represented by a finite poset. A category
with finitely many arrows has a nerve with a finite $2$-skeleton, so each
component of its classifying space has finitely presented fundamental
group. Its nerve need not be homotopy equivalent to a finite CW complex.
In particular the one-object category $BC_2$, with arrows $1,g$ and
$g^2=1$, is finite, whereas $B(BC_2)$ is not a compact object of $\Spc$.
\end{claim}
\begin{proof}
A finite CW complex is homotopy equivalent to a finite simplicial complex
$L$ by simplicial approximation and the mapping-cylinder construction;
we use precisely the finite clause of \citet[Theorem~2C.5]{Hatcher2002}.
For completeness, the inductive construction replaces each attaching map
from a triangulated sphere by a simplicial approximation, and attaches a
finite simplicial mapping cone. The homotopy between the attaching maps
induces the comparison equivalence of their mapping cones. Finitely many
cells require finitely many such replacements. The poset $P(L)$ of
nonempty faces has nerve the barycentric subdivision of $L$, whose
realization is homeomorphic to $|L|$. The empty space is represented by
the empty poset.

For a finite category $A$, its nondegenerate $1$-simplices are its
nonidentity arrows and its nondegenerate $2$-simplices are composable
pairs of nonidentity arrows. Both sets are finite. In any connected
component, choose a maximal tree in the finite $1$-skeleton. The remaining
edges are a finite set of fundamental-group generators and the finitely
many $2$-simplices give a finite set of relations; cells of dimension
at least three do not change the fundamental group. This also proves
finiteness of the number of components.

In the normalized mod-$2$ bar complex of $BC_2$ there is exactly one
nondegenerate simplex $(g,\ldots,g)$ in every degree. Every internal face
is degenerate, because it replaces $g,g$ by $1$, and the two end faces
cancel modulo $2$. Consequently $H^q(B(BC_2);\mathbf F_2)\cong\mathbf F_2$
for every $q\geq0$. This rules out finite CW homotopy type. More precisely,
write $B(BC_2)\simeq\mathbf {RP}^{\infty}$ and use the filtered system
$\mathbf {RP}^{r}\hookrightarrow\mathbf {RP}^{r+1}$. If its identity
factored in $\Hot$ through any $\mathbf {RP}^{r}$, its pullback on
$H^q(-;\mathbf F_2)$ for $q>r$ would be zero, contrary to the identity.
Thus the canonical map
\[
 \colim_r[\mathbf {RP}^{\infty},\mathbf {RP}^{r}]
 \longrightarrow
 [\mathbf {RP}^{\infty},\mathbf {RP}^{\infty}]
\]
is not surjective. This proves noncompactness without replacing a
specified comparison by an abstract bijection. The multiplication table,
all associativity equations, and normalized boundaries through degree
nine are checked in \texttt{code/f\_checks.py}; the preceding symbolic
calculation proves the assertion in every degree. The cohomology-ring
calculation is also \citet[Theorem~3.19]{Hatcher2002}.
\end{proof}

\subsection{Set-valued tests}

Put
\[
 R_1^{\mathrm{cat}}:\Hot\longrightarrow
 \Fun((h\mathcal A_{\mathrm{cat}})^{\op},\Set),
 \qquad R_1^{\mathrm{cat}}(Y)(A)=[A,Y].
\]
The same notation with superscript $\mathrm{fin}$ uses
$h\mathcal A_{\mathrm{fin}}$.

\begin{claim}{F.2}{Refuted}{Faithfulness and conservativity of finite set-valued tests}
Both $R_1^{\mathrm{cat}}$ and $R_1^{\mathrm{fin}}$ fail to be faithful
and fail to be conservative. There is a single endomorphism
$s_B:BG\to BG$ such that $s_B\ne\id_{BG}$, $s_B$ is not an isomorphism
in $\Hot$, and both restricted functors send it to the identity.
\end{claim}
\begin{proof}
Let $G$ be the group of permutations of $\N=\{0,1,\ldots\}$ with finite
support. Define an injective group homomorphism $s:G\to G$ by
\[
 s(\sigma)(0)=0,\qquad s(\sigma)(i+1)=\sigma(i)+1.
\]
It is a homomorphism by substitution in the displayed formula. It is not
surjective because its image fixes $0$, whereas $(0\ 1)\in G$ does not.
Let $s_B=Bs$ be the map of classifying spaces, hence a morphism of
$\Hot$. It cannot be an isomorphism because its map on the fundamental
group is $s$.

If $P$ is any finitely generated group and $\rho:P\to G$ a homomorphism,
the supports of the images of finitely many generators have a finite
union. Choose $m$ such that this union lies in $\{0,\ldots,m\}$. Every
$\rho(p)$ fixes the complement of this set. For
$c=(0\ 1\ \cdots\ m+1)$, considered as an element of $G$, direct
conjugation gives
\[
       c\rho(p)c^{-1}=s(\rho(p))\qquad(p\in P).
\]
The program \texttt{f\_checks.py} verifies this identity for all
$5913$ permutations supported on initial sets of sizes one through
seven; the formula itself proves the identity for arbitrary finite
support.

For a connected CW complex $A$, the canonical classification of maps
to a $1$-type gives
\[
 [A,BG]\cong\Hom(\pi_1(A,a),G)/G,
\]
where $G$ acts by conjugation. One obtains this classification by
constructing the map on the $1$-skeleton from a homomorphism, extending
across $2$-cells using its relations, and then extending uniquely up to
homotopy over higher cells because $\pi_i(BG)=0$ for $i>1$.
Changing the image-basepoint path is exactly conjugation. Thus the
preceding calculation says $(s_B)_*=\id$ on $[A,BG]$ whenever
$\pi_1(A,a)$ is finitely generated. Apply F.1 to every connected component
of a finite-category nerve, or of a finite CW complex. Maps from a
disjoint union are the product of maps from its components, and the
empty test gives the singleton set. Hence the induced natural
endomorphism of either restricted presheaf is the identity. It agrees
with the image of $\id_{BG}$ although the two morphisms of $\Hot$ are
distinct. This proves both failures.

This is the finite-support instance attributed to
\citet[Proposition~4.1]{MatumotoMinamiSugawara1984} by
\citet[Remark~2.3]{ArlinChristensen2023}. We have checked the latter source;
the former paper was not accessible in this run, so its bibliography
entry is explicitly unverified. The proof above does not depend on
that attribution.
\end{proof}

\begin{claim}{F.3}{Refuted}{Fullness of finite set-valued tests}
The functors $R_1^{\mathrm{cat}}$ and $R_1^{\mathrm{fin}}$ are not full.
This is a separate failure from F.2: there are natural transformations
between their values which are not induced by any morphism in $\Hot$.
\end{claim}
\begin{proof}
Choose an uncountable $\U$-small set $I$ and let $F_I$ be the free group
with basis $\{x_i:i\in I\}$. Put $X=BF_I$ and keep the countable group
$G$ from F.2. For each finite subset $J\subset I$, choose pairwise
disjoint transpositions $t_j\in G$, $j\in J$, and let
$\lambda_J:F_J\to G$ be the homomorphism with
$\lambda_J(x_j)=t_j$. Two choices for the same finite labelled set are
simultaneously conjugate in $G$: match the finitely many pairs in their
supports and extend the resulting bijection by a permutation with finite
support. The same observation says that, if $J\subset J'$, the
restriction of $\lambda_{J'}$ to $F_J$ is conjugate to $\lambda_J$.

For a connected test $A$ in either test category, a homomorphism
$\rho:\pi_1(A,a)\to F_I$ has finitely generated image, and therefore
its image lies in $F_J$ for a finite set of basis letters $J$. Define
\[
 \eta_A:[A,X]\longrightarrow[A,BG],
 \qquad [\rho]\longmapsto[\lambda_J\rho].
\]
Passing to $J\cup J'$ proves independence of the chosen finite support.
If $\rho$ is replaced by a conjugate $h\rho h^{-1}$ in $F_I$, enlarge
$J$ to contain the finitely many letters of the word $h$; then applying
$\lambda_J$ changes the image by conjugation by $\lambda_J(h)$.
Thus $\eta_A$ is defined on free homotopy classes. Precomposing a
group homomorphism with the homomorphism induced by a map of test
spaces commutes with this construction; changes of basepoint are the
same conjugacies already checked. On disconnected test spaces use the
construction on each component. It respects maps which send different
components to the same component, because each component is treated
by precisely the same precomposition rule. For the empty test there
is the unique function between singleton sets. We have therefore
constructed an actual natural transformation, not merely a family of
functions.

Suppose $\eta$ were represented by $f:X\to BG$. Choose a group
homomorphism $\lambda:F_I\to G$ representing $f$. Testing the loop
$x_i$ by $S^1$ says that $\lambda(x_i)$ is a transposition. Testing
the pair $x_i,x_j$ by $S^1\vee S^1$ says that the pair
$(\lambda(x_i),\lambda(x_j))$ is simultaneously conjugate to two
disjoint transpositions. In particular $\lambda(x_i)\ne\lambda(x_j)$
when $i\ne j$. Both test spaces have finite poset presentations by
F.1. Thus $i\mapsto\lambda(x_i)$ injects the uncountable set $I$ into
the countable set $G$, a contradiction. All sets and classifying
spaces used here belong to the fixed universe. We searched the primary
sources listed in the report for this particular fullness argument;
we do not attach a priority claim to it.
\end{proof}

\begin{claim}{F.4}{Known}{No small set of unpointed tests repairs conservativity}
For every $\U$-small set $\mathcal G$ of $\U$-small spaces, there is
an endomorphism $f:X\to X$ of a $\U$-small connected $1$-type which
is not an isomorphism in $\Hot$ and satisfies
$[S,f]=\id_{[S,X]}$ for every $S\in\mathcal G$. In particular no such
set jointly reflects isomorphisms or separates all morphisms.
\end{claim}
\begin{proof}
This is the construction proving
\citet[Theorem~2.1]{ArlinChristensen2023}; we give its size argument.
Choose an uncountable regular cardinal $\alpha\in\U$ strictly larger
than the cardinality of $\pi_1(S,s)$ for every $S\in\mathcal G$ and
every component basepoint $s$. The collection of these groups is a
$\U$-small set, so a successor of a sufficiently large cardinal in
$\U$ has this property. Regard $\alpha$ as its initial ordinal and put
\[
 H=\bigcup_{\beta<\alpha}\operatorname{Sym}(\beta)
       \subset\operatorname{Sym}(\alpha),
\]
where each permutation fixes all ordinals outside $\beta$. This is a
$\U$-small group. Let $j:\alpha\hookrightarrow\alpha$ send $\gamma$
to $\gamma+1$. Define $s_H:H\to H$ by transporting a permutation along
$j$ on the successor ordinals and fixing every limit ordinal, including
$0$. Products are preserved, the map is injective, and it is not
surjective because every image fixes $0$. A bounded support remains
bounded after this operation, so it does map $H$ into $H$.

If $|P|<\alpha$ and $\rho:P\to H$, regularity bounds the supports of
all $\rho(p)$ in a single infinite limit ordinal $\beta<\alpha$.
There is a permutation $\tau\in H$ with
$\tau(j(\gamma))=\gamma$ for all $\gamma<\beta$: extend this
bijection from $j(\beta)$ to $\beta$ over a sufficiently large
bounded interval of cardinality $|\beta|$, using an extra interval
of that cardinality to make the two complements equinumerous, and
fix its complement in $\alpha$. The bound remains below $\alpha$
because its cardinality is below $\alpha$. Since $s_H\rho(p)$ is
supported on $j(\beta)$, we obtain
$\tau s_H\rho(p)\tau^{-1}=\rho(p)$ on that support and on its
complement. Thus $Bs_H$ induces the identity on maps from every
connected component of every $S\in\mathcal G$. Products over the
components finish the proof, including arbitrarily many components.
\end{proof}

\begin{claim}{F.5}{Known}{The stronger nonconcreteness obstruction}
The category $\Hot$ admits no faithful functor to $\Set_{\U}$ with
the usual fixed-universe interpretation. The obstruction already
occurs on simply connected, finite-dimensional CW complexes. Therefore
faithfulness of a restricted set-valued Yoneda functor cannot be
obtained by substituting any other $\U$-small full test subcategory.
\end{claim}
\begin{proof}
We use the main theorem of \citet[p.~1 and the conclusion on p.~5 of
the 2004 reprint]{Freyd1970}. The theorem states the failure of every
homotopy-invariant set-valued functor to separate a nonnull map from
the null map. The construction on p.~5 works on $(r+3)$-dimensional,
$r$-connected CW complexes for every $r\geq1$, so based versus
unbased homotopy classes cause no exception. This is an invocation
of the proved theorem, not a new proof of nonconcreteness. For a
$\U$-small category $\mathcal A$, the functor
\[
 \Fun(\mathcal A^{\op},\Set_{\U})\longrightarrow\Set_{\U},
 \qquad F\longmapsto\coprod_{A\in\Ob\mathcal A}F(A)
\]
is faithful: the tag $A$ retains the component of each natural
transformation. A faithful restricted Yoneda functor followed by
this functor would contradict Freyd's theorem.
\end{proof}

\begin{claim}{F.6}{Proved}{The telescope phantom and the failure of epimorphy}
There is an essential unpointed map
$c:T\to K(\Z,2)$, where
$T=\operatorname{Tel}(S^1\xrightarrow{2}S^1\xrightarrow{2}\cdots)$,
such that $c\circ a$ is nullhomotopic for every map $a:BA\to T$
from a finite category. If $T_r$ denotes the $r$th finite
subtelescope, the canonical map $(T_r)_r\to T$ in $\Ind(\Hot)$ is
not an epimorphism.
\end{claim}
\begin{proof}
The telescope has fundamental group $\Z[1/2]$ by van Kampen. Any
map from a sphere of dimension at least two, and any specified
nullhomotopy, lies in a finite subtelescope after cellular
approximation; the latter retracts onto a circle. Hence
$T\simeq K(\Z[1/2],1)$. Contracting the ray of vertices gives a
homotopy equivalent CW complex with one vertex, loops $e_r$, and
one $2$-cell for each relation $e_r=2e_{r+1}$ in its cellular
boundary. Its cochain differential is
\[
 d:\prod_{r\geq0}\Z\longrightarrow\prod_{r\geq0}\Z,
 \qquad d((b_r))=(b_r-2b_{r+1})_r.
\]
Consequently $H^2(T;\Z)=\operatorname{coker}d$. Set $a_{2j}=1$
and $a_{2j+1}=0$. If $a=d(b)$, then for every $N\geq0$
\[
 b_0=\sum_{j=0}^{N}2^ja_j+2^{N+1}b_{N+1}.
\]
In the $2$-adic integers the right-hand partial sum tends to
$1+4+16+\cdots=-1/3$. Since $\Z\hookrightarrow\Z_2$ is injective,
this would give $3b_0=-1$ in $\Z$, which is impossible. Thus $[a]$
is a specified nonzero element of $H^2(T;\Z)$. Its representing
map $c$ exists by \citet[Theorem~4.57]{Hatcher2002}; that theorem
also gives the unpointed identification because the target is
simply connected. The entire cohomology group is
$\Z_2/\Z$, as calculated in
\citet[Examples~3F.7 and~3F.9]{Hatcher2002}, although nonvanishing
alone suffices here.

Every homomorphism from a finitely generated group into
$\Z[1/2]$ has cyclic image, since finitely many denominators have
a common power of two. F.1 and the classification of maps into a
$1$-type therefore imply that a map $BA\to T$, on each connected
component, factors in $\Hot$ through a circle or a point.
Because $[S^1,K(\Z,2)]=0$, its composite with $c$ is nullhomotopic.
No finite-dimensionality assertion about $BA$ is used. A finite
subtelescope also retracts onto a circle, whence
\[
 \Hom_{\Ind(\Hot)}((T_r)_r,K(\Z,2))
       =\lim_r[T_r,K(\Z,2)]=\{0\}.
\]
Precomposing the two distinct arrows $c,0:T\rightrightarrows K(\Z,2)$
with $(T_r)_r\to T$ therefore makes them equal. This is precisely
the failure of that arrow to be an epimorphism.
\end{proof}

\subsection{Space-valued tests}

\begin{claim}{F.7}{Proved}{The exact Ind-completion assertion}
In the fixed universe there is an equivalence
\[
 \Ind_{\omega}(\mathcal A_{\mathrm{fin}})\simeq\Spc
\]
whose restriction to $\mathcal A_{\mathrm{fin}}$ is its given
inclusion. Thus
\[
 R_{\infty}^{\mathrm{fin}}:\Spc\longrightarrow
 \Fun(\mathcal A_{\mathrm{fin}}^{\op},\Spc),
 \qquad X\longmapsto\Map(-,X)|_{\mathcal A_{\mathrm{fin}}}
\]
is fully faithful. The compact objects of $\Spc$ are exactly the
retracts in $\Spc$ of finite CW homotopy types; no assertion that
all such retracts have finite CW type is needed.
\end{claim}
\begin{proof}
The point is compact: its derived mapping-space functor is the
identity. Mapping out of a finite homotopy colimit gives a finite
homotopy limit, and filtered colimits of spaces commute with finite
homotopy limits \citep[Proposition~5.3.3.3]{Lurie2009}. The latter statement can be checked on Kan models:
a map from a finite simplicial set and any of its finitely many
specified horn fillings occurs at one stage; equality and homotopies
of such finite data occur at a later common stage. A finite CW
complex is obtained from the empty space and points by finitely
many homotopy pushouts attaching cells. Induction proves its
compactness for filtered colimits of spaces, not just compactness
of its set of homotopy classes.

A simplicial model for any $\U$-small space is the union of its
finite simplicial subsets. This union is filtered and is its
homotopy colimit, since inclusions are cofibrations; its finite
pieces have finite CW realizations. Hence finite CW homotopy types
generate $\Spc$ under filtered colimits. These two facts verify
all three hypotheses of
\citet[Proposition~5.3.5.11(2)]{Lurie2009}, giving the claimed
equivalence. Alternatively, if
$X=\colim_iK_i$ and $Y=\colim_jL_j$ are such filtered presentations,
compactness and the universal property of the source colimit give
\[
 \Map(X,Y)\simeq
 \lim_i\Map(K_i,Y)\simeq
 \lim_i\colim_j\Map(K_i,L_j),
\]
which is exactly the mapping-space formula in the Ind-category.
These are canonical comparison equivalences, natural in both
variables. The restricted Yoneda functor is the defining embedding
of that Ind-category into presheaves.

If $X$ is compact, compactness applied to its finite-piece
presentation makes $\id_X$ factor in $\Hot$ through some $K_i$.
The corresponding homotopy $r i\simeq\id_X$ gives a homotopy
retract; replacing $i$ by its mapping cylinder realizes the retract
datum in $\Spc$. Conversely a retract of a compact object is compact,
because its mapping-space functor is a retract of the corresponding
compact mapping-space functor. This is also the idempotent-completion
statement \citet[Lemma~5.4.2.4]{Lurie2009}.
\end{proof}

\begin{claim}{F.8}{Proved}{Finite-category space-valued tests are dense, but are not all compact}
The functor
\[
 R_{\infty}^{\mathrm{cat}}:\Spc\longrightarrow
 \Fun(\mathcal A_{\mathrm{cat}}^{\op},\Spc)
\]
is fully faithful. Nevertheless the filtered-colimit extension
$\Ind_{\omega}(\mathcal A_{\mathrm{cat}})\to\Spc$ of the inclusion
is not an equivalence. Thus ``spaces are the Ind-completion of
finite spaces'' proves the intended positive statement only after
specifying finite CW homotopy types, or their retract completion,
as the compact generators.
\end{claim}
\begin{proof}
F.1 gives an inclusion
$i:\mathcal A_{\mathrm{fin}}\hookrightarrow\mathcal A_{\mathrm{cat}}$.
For $X\in\Spc$, the presheaf $R_{\infty}^{\mathrm{cat}}X$ is the
right Kan extension of $R_{\infty}^{\mathrm{fin}}X$ along
$i^{\op}$. Indeed at $A\in\mathcal A_{\mathrm{cat}}$ the
right-Kan-extension formula is
\[
 \lim_{(K\to A)\in(\mathcal A_{\mathrm{fin}}/A)^{\op}}
       \Map(K,X),
\]
where $\mathcal A_{\mathrm{fin}}/A=\mathcal A_{\mathrm{fin}}
\times_{\Spc}\Spc_{/A}$. The density statement in F.7
identifies this canonical limit with $\Map(A,X)$. Restriction
followed by this right Kan extension is the identity on these
presheaves; right Kan extension along a fully faithful inclusion
is fully faithful. Therefore, using F.7, maps between
$R_{\infty}^{\mathrm{cat}}X$ and $R_{\infty}^{\mathrm{cat}}Y$
are exactly maps $X\to Y$, with their entire mapping spaces.

Every representable object of
$\Ind_{\omega}(\mathcal A_{\mathrm{cat}})$ is compact. An
equivalence extending the inclusion would consequently make
$B(BC_2)$ compact in $\Spc$, contrary to F.1. The negative
assertion concerns the specified extension of the inclusion;
we do not infer it merely from the fact that a nerve is
infinite-dimensional.
\end{proof}

\subsection{Finite categorical height}

For each integer $n\geq1$, let $h_n\Spc$ have the same objects as
$\Spc$ and mapping spaces $\tau_{\leq n-1}\Map(X,Y)$. We use
``$n$-category'' in this $(n,1)$ sense. This convention retains
homotopy classes of homotopies at $n=2$; it does not mean the
category of $n$-truncated spaces. For either test subcategory
$\mathcal A$, put
\[
 R_n^{\mathcal A}(X)=
 \tau_{\leq n-1}\Map(-,X)|_{\mathcal A}
 \ \in\ \mathcal P_{n-1}(\mathcal A)
 :=\Fun(\mathcal A^{\op},\Spc_{\leq n-1}).
\]
Because the target is an $n$-category this functor factors through
$h_n\Spc$. The associated functor on the homotopy category has
maps given by coherent natural transformations modulo coherent
homotopy. At $n=2$ these are pseudonatural transformations and
modifications in an equivalent groupoid-enriched model. In
particular, a collection of unrelated homotopies at individual
test spaces is not a modification.

\begin{claim}{F.9}{Known}{Conservativity returns at height two}
For every integer $n\geq2$, both $R_n^{\mathrm{fin}}$ and
$R_n^{\mathrm{cat}}$ reflect equivalences. In fact a map
$f:X\to Y$ is a weak equivalence if
\[
 \Pi_1\Map(S^r,f):\Pi_1\Map(S^r,X)
             \longrightarrow\Pi_1\Map(S^r,Y)
\]
is an equivalence of groupoids for every $r\geq0$.
\end{claim}
\begin{proof}
The second assertion is
\citet[Theorem~4.3]{ArlinChristensen2023}. We give a direct proof
which also checks disconnected and empty spaces. The point is a
retract of $S^0$, so the hypothesis implies that
$\Pi_1 f:\Pi_1X\to\Pi_1Y$ is an equivalence. Thus $f$ is a
bijection on components and an isomorphism on every based
fundamental group. If one space is empty, the other is empty and
there is nothing further to prove. Fix a component and a basepoint
$x\in X$, and use $f$ to identify the two fundamental groups with
a group $\Gamma$.

For $r\geq2$, free homotopy classes of maps from $S^r$ into this
component are the orbit set of $\pi_r(X,x)$ under the
$\Gamma$-action. This follows by recording the path traced by
the basepoint during a free homotopy. The map
$u_r:\pi_r(X,x)\to\pi_r(Y,f(x))$ is a
$\Gamma$-equivariant group homomorphism, and the hypothesis gives
a bijection of its orbit sets. Surjectivity of the orbit map
implies surjectivity of $u_r$: write a target element as
$\gamma u_r(a)=u_r(\gamma a)$. If $u_r(a)=0$, injectivity of the
orbit map forces the orbit of $a$ to equal the orbit of $0$,
which is the singleton $\{0\}$. Thus the kernel is zero and
$u_r$ is injective. This proves the isomorphism on every
homotopy group, hence the weak equivalence. For $n>2$ apply
$\tau_{\leq1}$ to the assumed equivalences of mapping
$(n-1)$-types, and use the case just proved. Spheres and the
point have finite poset presentations, so both test categories
contain all the tests used.
\end{proof}

\begin{claim}{F.10}{Refuted}{Faithfulness at the highest retained dimension}
For every finite $n\geq1$, neither $R_n^{\mathrm{fin}}$ nor
$R_n^{\mathrm{cat}}$ is faithful on the highest retained
morphisms. More precisely, for $T$ from F.6, $Y_n=K(\Z,n+1)$,
and the constant map $0:T\to Y_n$, the comparison on
$\pi_{n-1}$ of mapping spaces has a nonzero kernel. For $n=1$
this means the map on components. Consequently these functors
are not fully faithful as $n$-functors for any finite $n$.
\end{claim}
\begin{proof}
Let $\alpha\ne0$ be the class in $H^2(T;\Z)$ constructed in F.6.
For $n\geq2$, the exponential adjunction and the Eilenberg--Mac~Lane
classification give
\[
 \pi_{n-1}\Map(T,K(\Z,n+1),0)
   \cong H^{n+1}(S^{n-1}\times T,\{*\}\times T;\Z)
   \cong H^2(T;\Z).
\]
The second isomorphism is the relative suspension isomorphism;
its integral sphere factor is free, so no Tor term occurs. The
first classifies families based at the constant map. At $n=1$
use $[T,K(\Z,2)]=H^2(T;\Z)$ instead. Pulling back this class
along any $a:A\to T$ from either finite-test category gives
$a^*\alpha=0$ by F.6. Hence the associated top-dimensional
morphism becomes zero at every object and every mapping-space
component of the test presheaf.

A top-dimensional morphism between maps of $(n-1)$-truncated
presheaves is determined by these components. One can see this
without a strict-model convention: its $(n-1)$-fold loop space
is a limit of discrete spaces, so evaluation at all objects and
all mapping components is injective on its set of points.
Therefore the image of $\alpha$ is zero as a coherent
presheaf morphism, not just at a selected basepoint. Since
$\alpha$ was nonzero before applying $R_n$, the asserted kernel
is nonzero. For $n=2$ this is an explicit phantom homotopy,
namely a nonidentity automorphism of the constant map in the
fundamental groupoid of $\Map(T,K(\Z,3))$.
\end{proof}

\begin{claim}{F.11}{Proved}{A precise gain for countable sources}
Let $X$ have countable CW homotopy type, let $Y$ be any $\U$-small
space, let $n\geq1$, and use finite CW tests. The canonical map
\[
 \theta_n:\tau_{\leq n-1}\Map(X,Y)
   \longrightarrow
 \Map_{\mathcal P_{n-1}(\mathcal A_{\mathrm{fin}})}
       (R_n^{\mathrm{fin}}X,R_n^{\mathrm{fin}}Y)
\]
is surjective on components. It induces isomorphisms on
$\pi_k$ for $0\leq k<n-1$ and surjections on $\pi_{n-1}$,
with basepoints wherever required. In particular coherent
height-two finite CW tests determine equality of homotopy
classes of maps with countable CW source. The same conclusion
about equality holds if one is given coherent tests on all
finite-category homotopy types.
\end{claim}
\begin{proof}
Choose finite subcomplexes
$K_0\subset K_1\subset\cdots$ exhausting a countable CW model
for $X$. Such an exhaustion exists by closure-finiteness and a
countable enumeration of its cells. These inclusions are
cofibrations, so the union is the homotopy colimit and
\[
 \Map(X,Y)\simeq\lim_r\Map(K_r,Y).
\]
Compactness of each test $A\in\mathcal A_{\mathrm{fin}}$ gives
$R_{\infty}^{\mathrm{fin}}X\simeq
\colim_r R_{\infty}^{\mathrm{fin}}K_r$ in the presheaf
$\infty$-category. Objectwise truncation is left adjoint to
inclusion of truncated presheaves. Combining that adjunction
with the preceding colimit and Yoneda gives a canonical
equivalence
\[
 \Map_{\mathcal P_{n-1}}(R_nX,R_nY)
      \simeq\lim_r\tau_{\leq n-1}\Map(K_r,Y).
\]
Under these identifications $\theta_n$ is precisely the
comparison between truncating a homotopy inverse limit and
taking the homotopy inverse limit of the truncations.

Write $m=n-1$. Each map
$\Map(K_r,Y)\to\tau_{\leq m}\Map(K_r,Y)$ has
$m$-connected homotopy fibres. Over any point of the
homotopy limit of the targets, the homotopy fibre of the
comparison is a homotopy inverse limit of a tower of
$m$-connected nonempty spaces. A tower of nonempty connected
spaces has nonempty homotopy inverse limit: after replacing
bonding maps by fibrations, choose points and their path lifts
recursively. More generally, the homotopy inverse limit of a
tower of $m$-connected spaces is $(m-1)$-connected. This
follows either by the same recursive extension over spheres,
or from the homotopy-limit fibre description as the homotopy
equalizer of the identity and shift on the product, whose long
exact sequence has zero homotopy groups in degrees below $m$.
For $m=0$ the assertion is only nonemptiness. Thus the
comparison is $m$-connected, which yields exactly the
surjectivity and isomorphisms stated. Truncating its source
at $m$ does not change these groups.

For $n=2$ the map on $\pi_0$ is a bijection. Thus a coherent
natural equivalence between the two images of maps
$f,g:X\to Y$ implies $f=g$ in $\Hot$. If the coherent
equivalence was given on $\mathcal A_{\mathrm{cat}}$, restrict
it to $\mathcal A_{\mathrm{fin}}$ before applying this
argument. We do not infer a fullness assertion on
$\mathcal A_{\mathrm{cat}}$, since its noncompact tests do not
satisfy the colimit comparison used above. F.10 shows that the
possible failure of injectivity on the highest group is real,
even for this countable source.
\end{proof}

\begin{claim}{F.12}{Refuted}{Ordinary fullness and faithfulness at every finite height}
For every finite $n\geq2$, neither $hR_n^{\mathrm{fin}}$ nor
$hR_n^{\mathrm{cat}}$ is full or faithful on the ordinary sets
of homotopy classes of maps. Nonfaithfulness has examples whose
source has CW cardinality $\aleph_{n-1}$; nonfullness has examples
whose source has CW cardinality $\aleph_n$. Moreover, for each
$0\leq k\leq n-1$, the map on the retained $\pi_k$ of a mapping
space can fail to be injective. All the groups, modules, and spaces
in these examples are $\U$-small.
\end{claim}
\begin{proof}
We first record the group-cohomological input and its cardinal control.
For $d\geq1$, put
\[
 G_d=\bigoplus_{\aleph_{d-1}}\Z/2.
\]
This is a locally finite abelian group of cardinality $\aleph_{d-1}$.
The established theorem
\citet[Corollary~6.10(2)]{BC_DKLT2002}, attributed there to Chen,
states that a torsion abelian group $G$ of cardinality $\aleph_r$
has cohomological dimension $r+1$ over any nonzero ring in which
the orders of its finite subgroups are invertible. In particular
\[
       \operatorname{cd}_{\Q}G_d=d.
\]
Thus there is a $\Q G_d$-module $M$ with
$H^d(G_d;M)\ne0$. We can choose such an $M$ of cardinality at
most $\aleph_{d-1}$: in the rational free bar resolution of the
trivial module $\Q$, let $\Omega^i\Q$ be its $i$th syzygy and
put $M=\Omega^d\Q$. The sequence
\[
 0\longrightarrow\Omega^d\Q\longrightarrow P_{d-1}
   \longrightarrow\Omega^{d-1}\Q\longrightarrow0
\]
does not split, since $\Omega^{d-1}\Q$ is not projective when
$\operatorname{pd}_{\Q G_d}\Q=d$. Its nonzero extension class
corresponds, by dimension shifting, to a nonzero element of
$\operatorname{Ext}^d_{\Q G_d}(\Q,M)=H^d(G_d;M)$.
Each bar module has cardinality $\aleph_{d-1}$, so the syzygies
have the required bound. For $d=2$ one may more explicitly take
$M=\Q G_2$, since \citet[Theorem~5.4]{BC_DKLT2002} proves that
$H^2(G_2;\Q G_2)\ne0$. These are invocations of the cited proved
group-cohomology theorems, not of Holt's conjectures for general
locally finite groups.

For a $\Q G$-module $M$ and $r\geq2$, write
\[
 L_G(M,r)=EG\times_G K(M,r)\longrightarrow BG
\]
for the split Eilenberg--Mac~Lane fibration, where the action on
$M$ is realized functorially on the simplicial abelian group
$K(M,r)$. It has the zero section. Its derived space of sections
over $BH\to BG$, for any subgroup $H\leq G$, is
\[
 \Gamma_H(M,r)=\Map_H(EH,K(M,r)),\qquad
 \pi_i\Gamma_H(M,r)=H^{r-i}(H;M)\quad(0\leq i\leq r),
 \tag{F.12a}
\]
and its homotopy groups above degree $r$ vanish. The equation for
$i=0$ is the classification of components as an abelian group.
One explicit verification uses the free simplicial $\Q H$-module
on $EH$: its normalized complex is the rational bar resolution.
Equivariant maps to the Dold--Kan model of $M[r]$ give the good
nonnegative truncation of the cochain complex
$\Hom_{\Q H}(P_\bullet,M)$ shifted by $r$. Its homology in
degree $i$ is precisely $H^{r-i}(H;M)$. This proves (F.12a),
including its naturality under restriction of subgroups.
When $H$ is finite, the invariants functor on $\Q H$-modules is
exact: an arbitrary lift of an invariant vector can be averaged
by $|H|^{-1}\sum_{h\in H}h$. Consequently
\[
 H^j(H;M)=0\ (j>0),\qquad
       \Gamma_H(M,r)\simeq K(M^H,r).
 \tag{F.12b}
\]

We next explain why this calculation controls coherent natural
transformations, rather than merely their values at individual tests.
Fix $n\geq2$, put $m=n-1$, and abbreviate
$\mathcal P_m(\mathcal A_{\mathrm{cat}})$ to $\mathcal P_m$.
For a locally finite group $G$, let $\mathcal F(G)$ be the filtered
poset of its finite subgroups. There is a canonical equivalence
\[
 R_n^{\mathrm{cat}}(BG)
     \simeq\colim_{H\in\mathcal F(G)}R_n^{\mathrm{cat}}(BH).
 \tag{F.12c}
\]
To verify this, evaluate first before truncation at a finite-category
test $A$. Each component of $A$ has finitely generated fundamental
group by F.1. Maps from that component to $BG$ form the groupoid
of group homomorphisms to $G$ and conjugating elements in $G$.
A homomorphism has finite image, since $G$ is locally finite,
and finitely many specified conjugating elements lie in a common
finite subgroup with that image. Thus this mapping groupoid is
the filtered colimit of the corresponding mapping groupoids for
$BH$. There are only finitely many components of $A$, so the
same assertion holds for the product over its components.
Taking nerves verifies the assertion on mapping spaces. Objectwise
truncation preserves colimits into $\mathcal P_m$, giving (F.12c).

Let $p:E\to BG$ be a fibration and consider the space of sections
of its image in $\mathcal P_m$. By (F.12c) and the universal
property of a colimit in the slice category, it is
\[
 \begin{aligned}
 &\Map_{\mathcal P_m/R_nBG}(R_nBG,R_nE)\\
 &\qquad\simeq
 \lim_{H\in\mathcal F(G)^{\op}}
 \operatorname{hofib}_{BH\hookrightarrow BG}
 \left(\tau_{\leq m}\Map(BH,E)
          \longrightarrow\Map(BH,BG)\right).
 \end{aligned}
 \tag{F.12d}
\]
Here $BH$ is itself a finite-category test, so truncated Yoneda
identifies the first mapping space inside each fibre exactly.
The second mapping space is a $1$-type, already $m$-truncated.
For a map with $m$-truncated base, taking the homotopy fibre
commutes with $m$-truncation of its total space. This follows
from the homotopy-group exact sequence, since the base has no
homotopy groups above degree $m$. Hence each term of (F.12d)
is $\tau_{\leq m}$ of the actual space of sections of $p$ over
$BH$. In particular, if those section spaces have homotopy type
$K(M^H,r)$ with $r>m$, every term of (F.12d) is contractible,
and so is their limit. This last conclusion is a fully coherent
statement; no compatible choice of individual nullhomotopies
has been assumed.

For nonfaithfulness, take $G=G_n$ and a nonzero
$\alpha\in H^n(G;M)$ as above, and put $E=L_G(M,n)$.
The zero section $s_0$ and a section $s_\alpha$ representing
$\alpha$ are different components of the actual section space
by (F.12a). They also determine different \emph{unpointed} homotopy
classes $BG\to E$. Indeed, the homotopy fibre sequence
\[
 \Gamma_G(M,n)\longrightarrow\Map(BG,E)
       \xrightarrow{p_*}\Map(BG,BG)
\]
shows that two fibre components can merge in the total space
only if they belong to the same orbit of
$\pi_1(\Map(BG,BG),\id)$. Composition with $s_0$ is a section
of $p_*$, so every loop of the base lifts to a loop at $s_0$.
Its action fixes the zero component. That orbit is a singleton,
and cannot contain $\alpha\ne0$.
On the other hand, (F.12b)--(F.12d) make the space of sections
of $R_np$ contractible. The two induced transformations are
therefore homotopic in $\mathcal P_m$. This proves failure of
ordinary faithfulness for finite-category tests. Restricting
the coherent homotopy proves the same failure for finite CW
tests.

For the additional assertion about every retained homotopy group,
keep $G=G_n$ and replace $E$ by $L_G(M,n+k)$ for
$1\leq k\leq n-1$. Formula (F.12a) puts the nonzero group
$H^n(G;M)$ in $\pi_k$ of its section space at zero.
The base $\Map(BG,BG)$ is a $1$-type, so the fibre exact
sequence injects this group into $\pi_k\Map(BG,E)$; for
$k\geq2$ it is an isomorphism onto that homotopy group.
Formula (F.12d) again makes the presheaf section space
contractible, so this nonzero element maps to zero.

For nonfullness, take $G=G_{n+1}$ and choose
$0\ne\beta\in H^{n+1}(G;M)$ with the same cardinal control.
Let $z,s_\beta$ be the zero and $\beta$ sections of
$L_G(M,n+1)\to BG$. Define $E_\beta\to BG$ to be the
fibrewise path fibration whose fibre at $b$ consists of paths
in the fibre of $L_G(M,n+1)$ from $z(b)$ to $s_\beta(b)$.
Its fibres are $K(M,n)$, and a homotopy section would give a
path from $z$ to $s_\beta$ in $\Gamma_G(M,n+1)$, contradicting
$\beta\ne0$. Thus $E_\beta\to BG$ has no homotopy section.
For finite $H\leq G$, (F.12b) implies that the two restricted
sections are connected by a path, and their path space has
homotopy type $K(M^H,n)$. Formula (F.12d) therefore provides
a section
\[
 \eta:R_nBG\longrightarrow R_nE_\beta
 \quad\text{with}\quad R_np\circ\eta\simeq\id_{R_nBG}
\]
in $\mathcal P_m$.

This section cannot be induced, even up to coherent homotopy,
by a map $f:BG\to E_\beta$. The reason, valid for either test
category $\mathcal A$, is the exact comparison for any $1$-type
$Z$ and $n\geq2$:
\[
\begin{split}
 \Map_{\mathcal P_m(\mathcal A)}(R_nX,R_nZ)
 &\simeq
 \Map_{\Fun(\mathcal A^{\op},\Spc)}
                 (R_\infty X,R_\infty Z)\\
 &\simeq\Map(X,Z).
\end{split}
 \tag{F.12e}
\]
The first equivalence is the truncation adjunction, because
$R_\infty Z$ already takes values in $1$-types; the second
is F.7 or F.8. If $R_nf\simeq\eta$, (F.12e) applied to
$Z=BG$ would give $pf\simeq\id_{BG}$, an impossible homotopy
section. Thus $hR_n^{\mathrm{cat}}$ is not full. Restrict
$\eta$ to finite CW tests; the same argument using (F.12e)
shows that it has no lift there either. Finally, the bar and
Eilenberg--Mac~Lane constructions and the finite homotopy
pullbacks used above preserve the indicated infinite cardinal
bounds. The entire construction remains in $\U$.
\end{proof}

This counterexample identifies why the tower argument in F.11 cannot
be extended to arbitrary filtered indexing categories. The canonical
comparison for finite CW tests remains
\[
 \tau_{\leq n-1}\lim_{i\in I^{\op}}\Map(K_i,Y)
   \longrightarrow
 \lim_{i\in I^{\op}}\tau_{\leq n-1}\Map(K_i,Y),
\]
but F.12 supplies actual failures of both injectivity and surjectivity
on its components. At height two the nonfaithful source has cardinality
$\aleph_1$, so the countable-source restriction in F.11 cannot simply
be removed. The higher Brown representability theorem
\citep[Theorem~4.2.11]{NguyenRaptisSchrade2022} concerns functors
on the full $n$-category satisfying its exactness hypotheses; it
does not assert fullness of these restricted finite-test presheaves.
Likewise the conservative comparison theorems in
\citet[Theorems~35 and~37]{Arlin2020} do not assert that these
ordinary morphism sets are faithfully represented.

\begin{center}
\begin{tabular}{L{0.16\linewidth}L{0.18\linewidth}L{0.22\linewidth}L{0.30\linewidth}}
\hline
Height & Equivalence detection & Highest-dimensional faithfulness & Ordinary map-class fullness and faithfulness\\
\hline
$n=1$ & False (F.2) & False (F.2) & Both false (F.2--F.3)\\
$2\leq n<\infty$ & True (F.9) & False (F.10, F.12) & Both false in general (F.12); finite CW tests give both for countable sources (F.11)\\
$n=\infty$ & True & True & Fully faithful, including every coherence (F.7--F.8)\\
\hline
\end{tabular}
\end{center}

\subsection{The identities sought in Chapters VII and IX}

\begin{claim}{F.13}{Proved}{What testing can justify for composition and reversal}
Let $P$ be a space with a specified morphism
$m:P\times P\to P$ in $\Hot$, and let $T:P\to P$ be a
specified endomorphism. Equality of the two associativity
composites, or of $T^2$ and $\id_P$, does not follow in general
from equality after every finite set-valued test. For each fixed
finite height there are actual operations and endomorphisms for
which the corresponding coherent finite-test identities hold
but these identities in $\Hot$ fail. A coherent
homotopy of their space-valued finite-test natural
transformations does imply the corresponding equality in
$\Hot$. A coherent height-two finite-test homotopy suffices
when the source of the proposed equality has countable CW
homotopy type. For arbitrary sources no fixed finite height
supplies a general equality criterion: F.12 gives distinct
ordinary homotopy classes with coherently identical finite tests
at every prescribed finite height.
\end{claim}
\begin{proof}
The associativity composites are the two typed maps
\[
 m\circ(m\times\id_P),\quad m\circ(\id_P\times m):
                  P\times P\times P\longrightarrow P
\]
in $\Hot$, using its Cartesian product. The reversal identity
compares $T\circ T,\id_P:P\to P$.

For a literal counterexample to both implications, fix $n\geq2$
and use the split fibration $p:E=L_G(M,n)\to BG$ and
$0\ne\alpha\in H^n(G;M)$ in F.12. Translation by the section
$s_\alpha$ defines a fibrewise automorphism
$T=t_\alpha:E\to E$, with inverse $t_{-\alpha}$.
F.12 gives a coherent path $R_ns_\alpha\simeq R_ns_0$ over
$R_nBG$. The functor $R_n$ preserves the fibre products over
$BG$ used in fibrewise addition: at each test the base mapping
space is a $1$-type, and truncation at $n-1\geq1$ preserves
such pullbacks. Therefore $R_nT\simeq\id_{R_nE}$.
But $T^2s_0=s_{2\alpha}$ is not homotopic to $s_0$ by the
unpointed argument in F.12, since the rational vector space
$H^n(G;M)$ has $2\alpha\ne0$. Thus $T^2\ne\id_E$.
Now put $m=T\circ\operatorname{pr}_1:E\times E\to E$.
Its finite-test image is coherently homotopic to
$\operatorname{pr}_1$, an associative operation. Its actual
associators are $T^2\operatorname{pr}_1$ and
$T\operatorname{pr}_1$ on $E^3$. If they were homotopic,
fixing the other two coordinates and composing with $T^{-1}$
would give $T\simeq\id_E$, contrary to
$Ts_0=s_\alpha\not\simeq s_0$. These examples also give
height-one false positives by applying $\pi_0$ to the
height-two coherent comparison. They refute the proposed
general testing implications, without asserting a failure of
any particular geometrically defined path operation.

A coherent space-valued homotopy
is a path in the mapping space of presheaves between the two
induced natural transformations. F.7 or F.8 identifies that
path with a homotopy of the original maps, proving their
equality in $\Hot$. If the domain is countable CW, F.11
already gives the required injectivity on $\pi_0$ at height
two. An equivalence-detection theorem has a premise about
invertibility of a single morphism and a conclusion about its
invertibility; it cannot be substituted for this premise
about a path between two parallel morphisms.

For the applications at VII ms.~127, VII ms.~150, and IX
ms.~13, one can therefore use a geometrically constructed
associator or reversal homotopy directly, or retain the
coherent data required by the preceding tests. Merely checking
that each finite family of paths gives equal elements of a
set supplies no such coherent datum. In a strict Moore-path
model, concatenation is an actual associative operation
because lengths add and functions are joined on successive
intervals; reversal of a path twice is the original path.
Transferring those identities to a manuscript model still
requires the specified comparison equivalence and
compatibility of its composition and reversal maps. The
testing theorem neither constructs that comparison nor
makes its compatibility automatic.
\end{proof}
